\section{Related Work}\label{sec:related}

\paragraph{Open small-LM training reports and reproductions.}
Small language models increasingly ship with detailed training documentation. Qwen3 releases open weights down to 0.6B parameters \citep{yang2025qwen3}; SmolLM2 documents a data-centric, multi-stage recipe at 1.7B and releases its curated datasets \citep{allal2025smollm2}; OLMo~2 releases weights, data, code, recipes, and intermediate checkpoints \citep{olmo20242}; and MiniCPM develops small-model training strategies systematically \citep{hu2024minicpm}. Closest to our setting, IMU-1 trains a 430M-parameter model with a bundle of modernizations -- the NorMuon optimizer, QK-norm, per-head gating, value residuals, and LayerNorm scaling, among others -- and reports approaching the benchmark performance of models trained on far more data, with per-component ablations \citep{grigorev2026imu}. On the data side, FineWeb documents ablation-driven web curation \citep{penedo2024fineweb}, and DataComp-LM provides a controlled testbed in which model-based filtering emerges as the key lever \citep{li2024datacomp}; our data-composition experiment places a dclm-edu arm against a FineWeb-Edu control at matched tokens (Table~\ref{tab:data}). These reports evaluate recipes as bundles at production scale. We work in the opposite direction: reproduce one released model bit-exactly (Table~\ref{tab:repro}), then re-test which published components survive single-variable, iso-FLOP, multi-seed attribution on a different base, data pipeline, and budget (Figure~\ref{fig:deconfound}).

\paragraph{Optimizers and baseline-tuning discipline.}
NorMuon adds neuron-wise normalization to Muon's orthogonalized updates and reports efficiency gains over both Adam and Muon at the 1.1B scale \citep{li2025normuon}, with AdamW \citep{loshchilov2017decoupled} the default comparator. Because optimizer wins are routinely attributed to undertuned baselines, our NorMuon-vs-AdamW comparison (+0.4743 BPB on wikitext-2, 95\% CI [+0.4435, +0.5052], 3 seeds, paired, iso-FLOP at 42M tokens per arm, text-lm-v2; Table~\ref{tab:attrib}) carries a matched-configuration AdamW learning-rate sweep whose entire spread is roughly one tenth of the gap in Table~\ref{tab:attrib}, and we scope the result as an early-training optimization-speed signal at one architecture and one budget, not a claim about behavior at scale.

\paragraph{Compute-optimal scaling and matched-compute comparison.}
Chinchilla showed that at fixed compute, parameters and training tokens should scale together, and it compared models at an equalized compute budget \citep{hoffmann2022training}; we follow the FLOP-accounting conventions of \citet{chowdhery2022palm}. This study inherits the matched-compute discipline rather than extending it: every comparison varies exactly one factor with training FLOPs matched within 5\%, and all budgets sit far below compute-optimal (the base run is $\approx$2 tokens per parameter; Table~\ref{tab:arms}), so effects are reported as fixed-budget attributions, never scaling laws.

\paragraph{Mid-training, annealing, and long-context evaluation.}
MiniCPM's warmup-stable-decay schedule frames the decay phase as a vehicle for continued training and domain adaptation \citep{hu2024minicpm}; OLMo~2 applies a dedicated late-stage mixture during annealing \citep{olmo20242}, and SmolLM2 rebalances its mixture across stages \citep{allal2025smollm2}. Such reports motivate premium-data anneals but seldom isolate the data effect from the learning-rate decay that accompanies it. Our anneal therefore carries an iso-token control under the identical cooldown: the control barely moves off the base, while the premium mix improves code BPB by +0.2716 (95\% CI [+0.2641, +0.2792], 3 seeds, paired, 150.7M tokens per cell, text-lm-v2; Table~\ref{tab:midtrain}). For long context we adopt the passkey retrieval probe of \citet{mohtashami2023landmark} as a per-rung accuracy ladder with Wilson intervals (ecl-ladder-v1; Figure~\ref{fig:ecl-ladder}), consistent in spirit with RULER's finding that claimed context lengths overstate effective ones \citep{hsieh2024ruler}.

\paragraph{RLVR skepticism at small scale.}
Recent work questions how much reinforcement learning with verifiable rewards adds beyond the base model. \citet{yue2025does} find via large-$k$ pass@$k$ that base models match or exceed their RLVR-trained counterparts, concluding that RLVR sharpens what the base model already contains rather than expanding it; \citet{shao2025spurious} show that randomly assigned rewards produce large MATH-500 gains on Qwen2.5-Math models specifically and urge validation across model families, which motivates our mandatory random-reward control on a Qwen-lineage base; \citet{liu2025understanding} identify optimization biases in GRPO and propose the unbiased Dr.~GRPO variant we adopt; and DAPO documents the system-level techniques needed to make large-scale RLVR reproducible \citep{yu2025dapo}. VibeThinker-3B argues the optimistic case, reporting benchmark results at 3B parameters that it claims rival far larger systems \citep{xu2026vibethinker}. We contribute a pre-registered small-scale data point: Dr.~GRPO with DAPO-style group filtering on our 0.6B base beats neither its SFT floor nor the random-reward gate on GSM8K \citep{cobbe2021training} or MATH-500 levels 1--3 \citep{hendrycks2021measuring}, with training reward flat at $\approx$0.9\% (Table~\ref{tab:rlvr}, Figure~\ref{fig:rlvr}; math-acc-v1, n=1 seed, directional by design), scored under the pass@$k$ protocol of \citet{chen2021evaluating} with Wilson intervals.

\paragraph{Positioning.}
This paper introduces no new training technique. Its contribution is single-variable, iso-FLOP, multi-seed attribution of published techniques at reproduction scale on one machine, anchored by a bit-exact reproduction, together with negative results reported at the same standard as wins: SFT response-masking that does not separate from its iso-FLOP control (Table~\ref{tab:sft}), the pre-registered GRPO null (Table~\ref{tab:rlvr}), and a context-extension direction gated off by a step-0 diagnostic.
